% !TeX root = ../../main.tex
\section{Problem and Tasks}
\label{sec:mobiwac:problem}

Given a user's time-ordered check-in history, we predict two properties of the next visit. The first is its \emph{category}, one of seven fixed labels: Community, Entertainment, Food,
Nightlife, Outdoors, Shopping, and Travel. Istanbul's source collection maps its places onto the same seven labels. The second is its \emph{region}, the place's census tract (for Istanbul, the \emph{mahalle}, a municipal neighborhood). We do not predict the exact next place; both properties are easier to learn and, for most uses, enough. Region, unlike category, is a large label set: we frame next-region as classification over candidate regions, from 520 classes (Istanbul) to 8,501 (California; Table~\ref{tab:mobiwac:datasets}).

The motivation is practical: the category says what type of place comes next, hence what a user will want; the region says where, hence where to prepare. The mobility literature supports such preparation: city services have long been built on location-based social network traces~\cite{silva2019urbancomputing}. A recent analysis of tourist check-in
mobility finds that visits concentrate on a few hub places, argues that crowds and demand therefore concentrate, and ties this structure to
infrastructure planning and adaptive strategies~\cite{moura2025mobilityaware}.
That analysis reads past visits; we predict two properties of
the next one, and aggregated over users, those predictions point to where
demand is heading before it arrives. A census
tract is a neighborhood, not a radio cell. We therefore scope our claims to
neighborhood-level preparation: demand and load anticipation, caching content
ahead of time, and capacity planning. Cell association and handover are
radio-level decisions and remain out of scope. We build and evaluate no such
service here; Section~\ref{sec:mobiwac:discussion} quantifies what these predictions
would give such a service.
